% Fragment giao diện được input trong
% Sections/06-HienThucVaTrienKhaiHeThong.tex -- Hiện thực và triển khai hệ thống.

% Các ảnh giao diện được trình bày theo nhóm để giữ nội dung trong chương chính.
\newcommand{\frontendimage}[2]{%
    \IfFileExists{#1}{%
        \includegraphics[width=\linewidth,height=#2,keepaspectratio]{#1}%
    }{%
        \fbox{%
            \begin{minipage}[c][#2][c]{0.96\linewidth}
                \centering
                Ảnh giao diện sẽ được bổ sung tại:\\[0.3em]
                \texttt{\detokenize{#1}}
            \end{minipage}%
        }%
    }%
}

\newcommand{\frontendpair}[6]{%
    \begin{figure}[H]
        \centering
        \begin{subfigure}[t]{0.48\textwidth}
            \centering
            \frontendimage{#1}{0.24\textheight}
            \caption{#2}
            \label{#3}
        \end{subfigure}\hfill
        \begin{subfigure}[t]{0.48\textwidth}
            \centering
            \frontendimage{#4}{0.24\textheight}
            \caption{#5}
            \label{#6}
        \end{subfigure}
    \end{figure}%
}

\newcommand{\frontendtriple}[9]{%
    \begin{figure}[H]
        \centering
        \begin{subfigure}[t]{0.31\textwidth}
            \centering
            \frontendimage{#1}{0.18\textheight}
            \caption{#2}
            \label{#3}
        \end{subfigure}\hfill
        \begin{subfigure}[t]{0.31\textwidth}
            \centering
            \frontendimage{#4}{0.18\textheight}
            \caption{#5}
            \label{#6}
        \end{subfigure}\hfill
        \begin{subfigure}[t]{0.31\textwidth}
            \centering
            \frontendimage{#7}{0.18\textheight}
            \caption{#8}
            \label{#9}
        \end{subfigure}
    \end{figure}%
}

\newcommand{\frontendtwoone}[9]{%
    \begin{figure}[H]
        \centering
        \begin{subfigure}[t]{0.48\textwidth}
            \centering
            \frontendimage{#1}{0.22\textheight}
            \caption{#2}
            \label{#3}
        \end{subfigure}\hfill
        \begin{subfigure}[t]{0.48\textwidth}
            \centering
            \frontendimage{#4}{0.22\textheight}
            \caption{#5}
            \label{#6}
        \end{subfigure}

        \par\medskip
        \begin{subfigure}[t]{0.48\textwidth}
            \centering
            \frontendimage{#7}{0.22\textheight}
            \caption{#8}
            \label{#9}
        \end{subfigure}
    \end{figure}%
}

\newcommand{\frontendonetwo}[9]{%
    \begin{figure}[H]
        \centering
        \begin{subfigure}[t]{0.48\textwidth}
            \centering
            \frontendimage{#1}{0.22\textheight}
            \caption{#2}
            \label{#3}
        \end{subfigure}

        \par\medskip
        \begin{subfigure}[t]{0.48\textwidth}
            \centering
            \frontendimage{#4}{0.22\textheight}
            \caption{#5}
            \label{#6}
        \end{subfigure}\hfill
        \begin{subfigure}[t]{0.48\textwidth}
            \centering
            \frontendimage{#7}{0.22\textheight}
            \caption{#8}
            \label{#9}
        \end{subfigure}
    \end{figure}%
}

\newcommand{\frontendsingle}[3]{%
    \begin{figure}[H]
        \centering
        \begin{subfigure}[t]{0.62\textwidth}
            \centering
            \frontendimage{#1}{0.25\textheight}
            \caption{#2}
            \label{#3}
        \end{subfigure}
    \end{figure}%
}

\paragraph{Tổng quan hiện thực Frontend.}

Frontend được hiện thực bằng React và tổ chức theo bốn lớp trách nhiệm: định tuyến, trang chức năng,
component dùng lại và lớp API. Các lớp này phối hợp để hiển thị dữ liệu, tiếp nhận thao tác và gửi
request đến Backend; quyền truy cập và dữ liệu nghiệp vụ vẫn được xác minh ở phía máy chủ. Khu vực
quản trị sử dụng tuyến riêng, còn các chức năng khám phá sản phẩm, tư vấn AI và đọc bài viết được
đặt trong các tuyến công khai hoặc yêu cầu đăng nhập tùy theo nghiệp vụ.

\begin{table}[H]
    \centering
    \caption{Ánh xạ các nhóm giao diện Frontend với chức năng chính}
    \label{tab:frontend-interface-map}
    \begin{tabularx}{\textwidth}{|>{\RaggedRight\arraybackslash}p{0.22\textwidth}|>{\RaggedRight\arraybackslash}p{0.25\textwidth}|X|}
        \hline
        \textbf{Nhóm giao diện} & \textbf{Tuyến tiêu biểu} & \textbf{Chức năng hiện thực} \\
        \hline
        Khám phá sản phẩm & \texttt{/}, \texttt{/category} & Hiển thị danh mục, danh sách sản phẩm, tìm kiếm, lọc và điều hướng đến chi tiết sản phẩm. \\
        \hline
        Hỗ trợ lựa chọn & \texttt{/prodetail/:id}, \texttt{/compare}, \texttt{/wishlist} & Xem chi tiết, so sánh, lưu yêu thích và xem nhanh sản phẩm. \\
        \hline
        Trợ lý AI & \texttt{/tu-van}, \texttt{/ai} & Hội thoại nhiều lượt, gửi văn bản hoặc hình ảnh và hiển thị sản phẩm được đề xuất. \\
        \hline
        Mua hàng & \texttt{/cart}, \texttt{/checkout} & Quản lý giỏ hàng, nhập thông tin nhận hàng, chọn thanh toán và xác nhận đơn. \\
        \hline
        Tài khoản & \texttt{/login}, \texttt{/register}, \texttt{/info}, \texttt{/order-history} & Xác thực, cập nhật hồ sơ và theo dõi lịch sử đơn hàng. \\
        \hline
        Quản trị & \texttt{/admin} & Quản lý sản phẩm, khách hàng, đơn hàng, khuyến mãi, coupon và bài viết. \\
        \hline
    \end{tabularx}
\end{table}

\subsubsection{Giao diện dành cho khách hàng}

\paragraph{Trang chủ và danh mục sản phẩm.}

Trang chủ là điểm bắt đầu của hành trình mua sắm, kết hợp thanh điều hướng, ô tìm kiếm, banner,
danh mục và các sản phẩm nổi bật. Từ trang chủ, người dùng có thể chuyển đến danh mục hoặc nhập
truy vấn để tìm sản phẩm. Trang danh mục hiển thị danh sách theo nhóm, hỗ trợ lọc theo thuộc tính,
sắp xếp và phân trang; các thao tác này được giữ trong cùng một luồng để người dùng không mất ngữ
cảnh duyệt sản phẩm.

\frontendpair
{Images/Frontend/giao diện trang chủ của hệ thống.png}{Trang chủ của hệ thống}{fig:frontend-home}
{Images/Frontend/giao diện danh mục, tìm kiếm và lọc sản phẩm.png}{Danh mục, tìm kiếm và lọc sản phẩm}{fig:frontend-category}

\paragraph{Chi tiết sản phẩm và hỗ trợ lựa chọn.}

Trang chi tiết tập trung các thông tin cần thiết trước khi mua gồm hình ảnh, tên, giá, biến thể,
tồn kho, mô tả và đánh giá. Người dùng có thể chọn biến thể, theo dõi trạng thái còn hàng, thêm sản
phẩm vào giỏ hoặc chuyển sang danh sách yêu thích. Chức năng xem nhanh và so sánh giúp đối chiếu
nhiều lựa chọn mà không phải rời khỏi ngữ cảnh danh mục.

\frontendpair
{Images/Frontend/giao diện chi tiết sản phẩm.png}{Chi tiết sản phẩm}{fig:frontend-product-detail}
{Images/Frontend/giao diện trang sản phẩm yêu thích.png}{Danh sách sản phẩm yêu thích}{fig:frontend-wishlist}

\paragraph{Tìm kiếm và tư vấn sản phẩm bằng AI.}

Chatbot nổi hỗ trợ đặt câu hỏi nhanh trong quá trình duyệt, còn trang trợ lý chuyên biệt cung cấp
không gian cho hội thoại nhiều lượt. Người dùng có thể gửi câu hỏi bằng văn bản hoặc ảnh tham chiếu;
Frontend hiển thị trạng thái gửi, đang xử lý, kết quả và cảnh báo khi dịch vụ AI không sẵn sàng.
Các sản phẩm được đề xuất nằm trong ngữ cảnh câu trả lời, cho phép mở trang chi tiết và tiếp tục
mua sắm. Giá và tồn kho hiển thị trên giao diện vẫn được lấy lại từ Backend.

\frontendpair
{Images/Frontend/giao diện chatbot.png}{Chatbot hỗ trợ mua sắm}{fig:frontend-chatbot}
{Images/Frontend/giao diện trợ lý mua sắm đa phương tiện.png}{Trợ lý mua sắm đa phương thức}{fig:frontend-ai-consulting}

\paragraph{Giỏ hàng và thanh toán.}

Giỏ hàng cho phép cập nhật số lượng, xóa mặt hàng và kiểm tra tổng tiền trước khi checkout. Tại
màn hình thanh toán, người dùng nhập thông tin nhận hàng, áp dụng mã giảm giá hợp lệ và kiểm tra
lại danh sách sản phẩm trước khi xác nhận. Sau khi Backend ghi nhận đơn COD thành công, giao diện
hiển thị kết quả và cung cấp đường dẫn đến lịch sử đơn hàng.

\frontendpair
{Images/Frontend/giao diện quản lý giỏ hàng.png}{Quản lý giỏ hàng}{fig:frontend-cart}
{Images/Frontend/giao diện thanh toán.png}{Checkout và lựa chọn phương thức thanh toán}{fig:frontend-checkout}

\frontendpair
{Images/Frontend/giao diện đặt hàng thành công.png}{Xác nhận đặt hàng thành công}{fig:frontend-order-success}
{Images/Frontend/giao diện lịch sử đơn hàng.png}{Lịch sử và trạng thái đơn hàng}{fig:frontend-order-history}

\paragraph{Đăng nhập, đăng ký và quản lý tài khoản.}

Giao diện đăng nhập cho phép khách hàng truy cập các chức năng cá nhân hóa; giao diện đăng ký kiểm
tra dữ liệu đầu vào trước khi tạo tài khoản. Sau khi đăng nhập, người dùng có thể cập nhật thông tin
cá nhân, quản lý địa chỉ và tiếp tục các thao tác mua sắm theo quyền được cấp. Các trạng thái tải,
lỗi xác thực và lỗi kết nối được hiển thị ngay tại biểu mẫu để người dùng biết cách xử lý.

\frontendtwoone
{Images/Frontend/giao diện đăng nhập.png}{Đăng nhập hệ thống}{fig:frontend-login}
{Images/Frontend/giao diện đăng ký.png}{Đăng ký tài khoản}{fig:frontend-register}
{Images/Frontend/giao diện tài khoản.png}{Quản lý thông tin tài khoản}{fig:frontend-account}

\subsubsection{Giao diện dành cho quản trị viên}

Khu vực quản trị sử dụng bố cục riêng và chỉ cho phép tài khoản có vai trò phù hợp truy cập. Thanh
điều hướng bên giúp chuyển giữa các nhóm nghiệp vụ; các màn hình danh sách dùng bảng, bộ lọc và
phân trang, trong khi thao tác tạo, sửa hoặc chuyển trạng thái được thực hiện qua biểu mẫu và hộp
thoại xác nhận.

\paragraph{Tổng quan quản trị và quản lý dữ liệu.}

Trang tổng quan cung cấp các chỉ số vận hành để quản trị viên nắm nhanh tình trạng sản phẩm, người
dùng và đơn hàng. Màn hình quản lý sản phẩm hỗ trợ theo dõi catalog và cập nhật thông tin hàng hóa;
màn hình quản lý người dùng phục vụ kiểm soát tài khoản, vai trò và trạng thái hoạt động.

\frontendonetwo
{Images/Frontend/giao diện tổng quan dành cho quản trị viên.png}{Tổng quan dành cho quản trị viên}{fig:frontend-admin-dashboard}
{Images/Frontend/giao diện quản lý sản phẩm.png}{Quản lý sản phẩm}{fig:frontend-admin-products}
{Images/Frontend/giao diện quản lý người dùng.png}{Quản lý người dùng}{fig:frontend-admin-users}

\paragraph{Quản lý đơn hàng.}

Màn hình đơn hàng cho phép tìm kiếm, xem chi tiết và chuyển trạng thái theo vòng đời đã thiết kế.
Thông tin khách hàng, mặt hàng, tổng tiền và lịch sử xử lý được trình bày cùng ngữ cảnh để hạn chế
thao tác nhầm. Việc thay đổi trạng thái được gửi đến Backend, nơi tiếp tục kiểm tra quyền và điều
kiện nghiệp vụ trước khi lưu.

\frontendsingle
{Images/Frontend/giao diện quản lý đơn hàng.png}{Quản lý đơn hàng}{fig:frontend-admin-orders}

\paragraph{Quản lý khuyến mãi và mã giảm giá.}

Màn hình khuyến mãi dùng để cấu hình chính sách giảm tự động theo thời gian hiệu lực và phạm vi áp
dụng. Mã giảm giá được quản lý riêng vì có vòng đời, điều kiện sử dụng và cách nhập khác với chương
trình khuyến mãi. Việc tách hai màn hình giúp quản trị viên kiểm soát các chương trình đang hoạt động
và giảm nhầm lẫn khi nhiều chính sách được áp dụng đồng thời.

\frontendpair
{Images/Frontend/giao diện quản lý khuyến mãi.png}{Quản lý chương trình khuyến mãi}{fig:frontend-admin-promotions}
{Images/Frontend/giao diện quản lý mã giảm giá.png}{Quản lý mã giảm giá}{fig:frontend-admin-coupons}

\paragraph{Quản lý bài viết.}

Giao diện bài viết hỗ trợ quản trị nội dung giới thiệu, hướng dẫn và tin tức của cửa hàng. Nội dung
được tách khỏi catalog để vòng đời biên tập không phụ thuộc vào thao tác cập nhật sản phẩm; biểu mẫu
cho phép tạo, sửa, xem trước và thay đổi trạng thái hiển thị theo quyền quản trị.

\frontendsingle
{Images/Frontend/giao diện quản lý bài viết.png}{Quản lý bài viết}{fig:frontend-admin-blogs}

\paragraph{Nguyên tắc phản hồi và khả năng sử dụng.}

Các giao diện thống nhất về phản hồi thao tác, trạng thái đang tải, trạng thái rỗng và thông báo lỗi.
Nút gửi được khóa trong thời gian xử lý để hạn chế request trùng; danh sách dài sử dụng phân trang
hoặc bộ lọc; thao tác ảnh hưởng đến đơn hàng và dữ liệu quản trị cần bước xác nhận. Việc tổ chức theo
component và kiểm tra trên nhiều kích thước màn hình giúp các biểu mẫu, bảng và hộp thoại vẫn dễ sử dụng.
